% Generated draft: requires mathematical and source review.
\documentclass[11pt]{article}
\usepackage[margin=25mm]{geometry}
\usepackage{fontspec}
\setmainfont{Latin Modern Roman}
\usepackage{amsmath,amssymb,mathtools,unicode-math}
\setmathfont{Latin Modern Math}
\usepackage{enumitem,needspace,xurl,hyperref}
\setlength{\parindent}{0pt}
\setlength{\parskip}{5pt}
\setlength{\emergencystretch}{3em}
\newcommand{\N}{\mathbb{N}}
\newcommand{\Z}{\mathbb{Z}}
\newcommand{\Q}{\mathbb{Q}}
\newcommand{\R}{\mathbb{R}}
\newcommand{\C}{\mathbb{C}}
\newcommand{\F}{\mathbb{F}}
\newcommand{\A}{\mathbb{A}}
\newcommand{\PP}{\mathbb{P}}
\newcommand{\E}{\mathbb{E}}
\newcommand{\eps}{\varepsilon}
\newcommand{\dd}{\,\mathrm{d}}
\begin{document}
\section*{UnsolvedMath ID: 30003426}
\textbf{Source problem code:} catalogue entry
\subsection{Definitions and mathematical statement}
Let \(W\) denote the Brownian motion driving a stock-price process \(S\), and let \(\sigma\) denote its volatility process. The source does not specify the dynamics of \(S\), the state space or regularity of \(\sigma\), or the pricing probability measure. For a common observation time set \(I\), write
\[
\mathcal G_{\sigma}=\sigma(\sigma_t:t\in I),\qquad
\mathcal G_W=\sigma(W_t:t\in I),
\]
where \(\sigma(\cdot)\) denotes the generated sigma-algebra. Independence of the two processes means
\[
\mathbb P(A\cap B)=\mathbb P(A)\mathbb P(B)
\quad\text{for every }A\in\mathcal G_{\sigma}\text{ and }B\in\mathcal G_W.
\]
Here \(\mathbb P\) denotes the probability measure under consideration; the source does not identify it or specify \(I\).

\textbf{Research question.} Determine how far this independence assumption can be relaxed by allowing dependence between \(\sigma\) and \(W\), while retaining tractable option pricing and permitting asymmetric implied-volatility smiles.

An implied volatility is a volatility parameter obtained by matching an option price to a specified benchmark pricing formula. A smile describes its variation with strike or a chosen moneyness coordinate at fixed maturity. The source supplies neither the benchmark nor the coordinate and center relative to which asymmetry is intended. It also gives no definition of tractability, admissible dependence class, option family, or criterion for comparing relaxations. Consequently, this is a research-program formulation: the supplied evidence does not determine a unique quantified theorem or characterization problem. These missing specifications must be resolved before a more precise mathematical target can be stated.
\subsection{Short English statement}
Investigate which forms of dependence between stock-price noise and volatility permit tractable option pricing and asymmetric implied-volatility smiles. The source leaves the model, pricing assumptions, tractability criterion, and meaning of smile asymmetry unspecified, so the draft retains these as unresolved qualifications.
\subsection{Sources}
\begin{itemize}
\item Local UnsolvedMath registry: supplied complete record with permanent ID \(30003426\), source problem code \texttt{catalogue entry}, title \emph{Dependent Volatility Drivers in Tractable Pricing Models}, in the Oberwolfach Reports Open Problems collection. No individual proposer is supplied; the proposed year is recorded as 2017.
\item Supplied report attribution: \emph{Mathematics of Quantitative Finance (2018)}, Oberwolfach Reports, DOI 10.4171/owr/2017/13, pp.~683--769. \url{https://doi.org/10.4171/owr/2017/13}. The background describes these reports as records of mathematical workshops and their open-problem discussions. No problem-specific page or speaker attribution is supplied.
\item Supplied literature reference: Alos--Leon, \emph{curvature of the smile in stochastic volatility models}. \url{https://doi.org/10.1137/16M1086315}.
\item Supplied literature reference: Alos--Le\'on--Vives, \emph{leverage effect in stochastic volatility}. \url{https://doi.org/10.1016/j.physa.2008.10.012}.
\item Supplied literature check date: \texttt{2026-08-21}. This is supplied dated curation, not an independently verified assessment. It labels the problem partially solved and reports relaxations in correlated, leverage, jump, and rough-volatility models, with Malliavin-calculus formulas for implied-volatility effects; it further states that a complete general theory for arbitrary dependence remains unavailable. The background explicitly calls this a dated literature triage derived from the supplied OWR curation and requests independent verification of time-sensitive status and cited claims.
\item The registry labels the statement \texttt{corrected\_verified} and says that notation, typography, and sentence structure were repaired while preserving mathematical content. This supplied editorial label does not constitute independent mathematical or source verification of this draft.
\end{itemize}
\end{document}
